\section{Conclusion}
\label{sec:conclusion}

We introduced a self-supervised framework for predicting mechanical grey swans - rare but physically plausible component failures that are critically underrepresented in training data. By combining JEPA pretraining (which captures waveform texture) with temporal contrastive learning (which captures degradation dynamics), and fusing with domain-informed spectral features, we achieve state-of-the-art RUL prediction: RMSE $0.055 \pm 0.004$, a $75.5\%$ improvement over time-only baselines.

Our mechanistic analysis reveals that different self-supervised objectives learn fundamentally different representations, explaining their complementary strengths: JEPA excels within-dataset while contrastive learning enables cross-dataset transfer ($+38\%$ vs.\ time-only, $p < 0.001$). This understanding motivates the hybrid architecture and provides guidance for practitioners choosing between pretraining strategies.

Several directions remain open. Extending the framework to multivariate sensor inputs with physics-aware masking could exploit cross-modal dependencies (e.g., vibration--current relationships). Validation on turbofan engines (C-MAPSS) and other mechanical domains would strengthen generalization claims. Finally, synthetic data augmentation using physics-informed degradation models could address the fundamental scarcity of rare failure trajectories that motivates this work.
